% Generated draft: requires mathematical and source review.
\documentclass[11pt]{article}
\usepackage[margin=25mm]{geometry}
\usepackage{fontspec}
\setmainfont{Latin Modern Roman}
\usepackage{amsmath,amssymb,mathtools,unicode-math}
\setmathfont{Latin Modern Math}
\usepackage{enumitem,needspace,xurl,hyperref}
\setlength{\parindent}{0pt}
\setlength{\parskip}{5pt}
\setlength{\emergencystretch}{3em}
\newcommand{\N}{\mathbb{N}}
\newcommand{\Z}{\mathbb{Z}}
\newcommand{\Q}{\mathbb{Q}}
\newcommand{\R}{\mathbb{R}}
\newcommand{\C}{\mathbb{C}}
\newcommand{\F}{\mathbb{F}}
\newcommand{\A}{\mathbb{A}}
\newcommand{\PP}{\mathbb{P}}
\newcommand{\E}{\mathbb{E}}
\newcommand{\eps}{\varepsilon}
\newcommand{\dd}{\,\mathrm{d}}
\begin{document}
\section*{UnsolvedMath ID: 30001712}
\textbf{Source problem code:} catalogue entry
\subsection{Definitions and mathematical statement}
Let \(\Omega\subset V\) be a quiver cycle in the setting intended by the source, and put \(d=\operatorname{codim}(\Omega;V)\). Write \(c^H_\mu(\Omega)\) for its cohomological quiver coefficients and \(c^K_\mu(\Omega)\) for its \(K\)-theoretic quiver coefficients; these superscripts distinguish the two uses of \(c_\mu(\Omega)\) in the supplied statement. Let \(I_H\) and \(I_K\) denote their respective intended index sets, with \(|\mu|\) the size used in the source convention. For a partition \(\lambda=(\lambda_1,\ldots,\lambda_r)\), its size is \(|\lambda|=\sum_{j=1}^r\lambda_j\); if \(\mu\) is a tuple of partitions, its total size is the sum of their sizes. The record does not specify which partition indices are intended.

The supplied conjecture asks whether, for every such quiver cycle, all three assertions hold:
\begin{enumerate}
\item Cohomological positivity:
\[
\forall\mu\in I_H,\qquad |\mu|=d\ \Longrightarrow\ c^H_\mu(\Omega)\geq 0.
\]
\item Finiteness of the \(K\)-theoretic expansion:
\[
\#\{\mu\in I_K:c^K_\mu(\Omega)\neq 0\}<\infty.
\]
\item Alternating signs under the stated singularity hypothesis:
\[
\Omega\text{ has rational singularities}
\ \Longrightarrow\
\forall\mu\in I_K,\quad
(-1)^{|\mu|-d}c^K_\mu(\Omega)\geq 0.
\]
\end{enumerate}
The rational-singularity hypothesis is attached only to the third assertion. The record does not define a quiver cycle, the ambient representation space \(V\), the coefficient expansions or their bases, or the index restrictions. It also does not specify the ground field, its characteristic, allowed quivers, or geometric conventions needed to interpret rational singularities. Accordingly, these formulas preserve the supplied conjecture as a statement schema whose precise domain and coefficient definitions require confirmation from the cited source.
\subsection{Short English statement}
For every quiver cycle in the intended source setting, the conjecture asks for nonnegative cohomological coefficients in codimension degree, finitely many nonzero \(K\)-theoretic coefficients, and alternating signs for those \(K\)-theoretic coefficients when the cycle has rational singularities. The supplied record leaves the underlying geometric setting and coefficient conventions unspecified.
\subsection{Sources}
\begin{itemize}
\item Local UnsolvedMath registry, permanent ID \(30001712\), source problem code \texttt{catalogue entry}, title: \emph{Positivity of Cohomological and K-Theoretic Quiver Coefficients}. Collection: Oberwolfach Reports Open Problems; proposed year recorded as 2011. No individual proposer is supplied.
\item \emph{Representation Theory of Quivers and Finite Dimensional Algebras} (2011), Oberwolfach Reports, pp.\ 523--608, DOI \texttt{10.4171/owr/2011/10}. Supplied report provenance and literature reference: \url{https://doi.org/10.4171/owr/2011/10}. No precise problem page or passage is supplied.
\item Buch, \emph{Quiver coefficients of Dynkin type}, supplied literature reference: \url{https://arxiv.org/abs/0708.3418}.
\item Buch--Kresch--Tamvakis, described in the record as \emph{quiver coefficients and Schubert constants}, supplied literature reference: \url{https://arxiv.org/abs/math/0311390}.
\item Literature assessment dated \texttt{2026-08-21}: the supplied curation labels the problem partially solved and reports cohomological positivity in major classes and finiteness for Dynkin-type quivers, while describing the arbitrary-quiver-cycle conjecture as not fully resolved. This is supplied dated curation, not an independently verified literature assessment. The background explicitly requests independent verification of time-sensitive status and cited claims.
\end{itemize}
\end{document}
